\scanpage{531}
\mspage{271}

\[
\Gamma_{\mathbb Q}\unsure{-|\longrightarrow}\Gamma_{1,1}
\longrightarrow \operatorname{Aut\,ext}\widehat{\mathcal T}_{1,1}^{+}
\longrightarrow \operatorname{Aut\,ext}\widehat{\mathcal T}_{1,1}^{+'},
\]
by the preceding theorem this composite is injective, hence also
$\Gamma_{\mathbb Q}\longrightarrow\Gamma_{1,1}$, q.e.d.

\uline{Remark} We really have the impression, with the proposition at the beginning, and
the preceding result christened ``theorem'', of having proved
the same thing twice, and with the same proof
at that! So it is time, after this detour, to make sure that this
really is so, i.e. that the homomorphism
$\widehat{\Pi}_{0,3}\longrightarrow\deleted{\unsure{\Pi}}\widehat{\mathcal T}_{1,1}^{+'}$ which we have just used
does indeed identify with the homomorphism
$\widehat{\Pi}_{0,3}\longrightarrow\widehat{\Pi}'_{0,3}$,
by means of a suitable isomorphism (which remains to be described),
commuting with the outer actions of $\Gamma$,
which we had set out to construct (cf. (10) to (20)) --- and we
lost it along the way, in

\scanpage{532}

trying to \deleted{\unsure{prove}}\inserted{\unsure{find}}
$\Pi'_{0,3}\xrightarrow{\sim}\mathcal T_{1,1}^{+'}$ by
factorization in the composite homomorphism
\[
\Pi_{0,3}\xrightarrow{\substack{\mathrm{can.}\\\unsure{\mathrm{surj}}}}
\Pi'_{0,3}\longrightarrow\mathcal T_{1,1}^{+'}
\]
--- and in a sense we went astray, for the ``natural'' homomorphism
$\unsure{\varphi|\rho}\,\Pi_{0,3}\longrightarrow\mathcal T_{1,1}^{+'}$
which we stumbled upon was \uline{not} \inserted{at all} the one we had
in mind: I was way off base. So I must return to the charge!

\scanpage{533}
\mspage{272}

\uline{37} Theory of Moduli of elliptic curves via Legendre
(rigidification of level 2). \unsure{)}

I realize that I have made some nevertheless quite simple things
rather complicated, in the preceding \S\. First a pang of conscience about surface topology:

\uline{Proposition} \unsure{Let|Let}\deleted{\illegible{1}}
$X$ be a paracompact topological\unsure{${}^{*}$} surface,
$G$ a discrete group operating properly on $X$, we assume that for every
$x\in X$, the stabilizer $G_x$ operates faithfully on a neighborhood of $x$
(so that $G$ operates faithfully --- in fact the two must be
equivalent) preserving the local orientation (hence $G_x$ is
a cyclic group), so that $Y=X/G$ is a paracompact topological surface.
Consider the set
\[
X^{!}=\{x\in X\mid G_x\ne(1)\}
\]
--- which is a discrete subset of $X$ stable under $G$ --- and the image
$Y^{!}\unsure{\subset|\in}Y$ of $X^{!}$ --- a discrete subset of $Y$ --- and
the ``signature'' on $(Y,Y^{!})$, for which the coefficient $d_y$ $(y\in Y^{!})$
is the ramification index $\operatorname{ord}G_x$ at the
$x\in X$ over $y$. For every $G$-covering $X'$ of $X$, let us
then consider $X'/G=Y'$ as a space over $Y$. Then:

\margin{\uline{Variant} $\mathcal G$ operates
\deleted{$X$ \illegible{2} $\mathcal G$}
\deleted{\unsure{operates}}\inserted{on} $X$, \unsure{with|its}
\uline{invariant} subgroup $G$ in $\mathcal G$
satisfying the conditions opposite. We then have, putting
\deleted{\illegible{1}} $Y=X/G$ with an \unsure{operation} of
$\mathcal G/G=\mathcal H$
on it:
\[
\deleted{\illegible{1}}\underline{\operatorname{Rev}}.
(X,\mathcal G)\unsure{\approxeq|\simeq}
\underline{\operatorname{Rev}}((Y,\underline d),\mathcal H),
\]
hence
\[
\Pi_1(X,\mathcal G,x)
\xrightarrow{\sim}
\Pi_1((Y,\underline d,\mathcal H),y).
\]
Application
\[
\mathfrak S_{0,3}\simeq
\mathcal T'_{1,1}
(\simeq\unsure{\mathrm{Gl}|\mathrm{Sl}}(2,\mathbb Z)/\pm1)
\]
extension of $\mathcal T_{0,3}=\mathfrak S_3
\times\mathbb Z/2$ by $\Pi_{03}$}
\scanpage{534}

a) $Y'$ is \deleted{a} \inserted{a topological surface,} a ramified covering of $Y$, subordinate to the signature $\underline d\,\deleted{\illegible{1}}=\underline d(X/Y)$.

b) The functor $X'\longmapsto Y'=X'/G$ is an equivalence of categories
\[
(1)\qquad \mathrm{Rev}(X,G)\overset{\approx}{\longrightarrow}\mathrm{Rev}_{\underline d}(Y)
\]

\uline{N.B.} We can make explicit a quasi-inverse functor, by associating to every ramified covering $Y'$ of $Y$ compatible with $\underline d$, \deleted{\illegible{2}} \inserted{the} ``normalized'' $(X,G)$-covering (in a topological sense easy to make explicit) of $X\times_Y Y'$. We have an analogous statement for regular locally noetherian schemes of dimension $1$, where the quasi-inverse functor is obtained by normalizing $X\times_Y Y'$.

\uline{Corollary} Let $x\in X-X^{!}$, $y$ its image in $Y$. We assume $Y$ connected (i.e. that $G$ operates transitively on the set of connected components of $X$). Then we have a canonical isomorphism
\[
(2)\qquad \Pi_1(X,G;\unsure{a|x})\overset{\sim}{\longrightarrow}\Pi_1((Y,\underline d),y)
\]

\scanpage{535}
\mspage{273}

whence, for $X$ connected, a canonical exact sequence
\[
(3)\qquad
1\longrightarrow
\deleted{\Pi_1((X,\unsure{G|1}),x)}\inserted{\Pi_1(X,x)}
\longrightarrow\Pi_1((Y,\underline d),y)\longrightarrow G\longrightarrow 1
\]

\uline{N.B.} We have \unsure{?} an analogous corollary in the scheme-theoretic context, $X$ and $Y$ being regular schemes of dimension $1$\unsure{.} --- or, as the case may be, smooth relative schemes of relative dimension $1$ over an $S$, but in this case one would have to (if I need it) take a little care to specify \unsure{the statement|the statements} in a reasonable form.

\uline{Example} Let us take $X=\mathbb P^1(\mathbb C)\setminus\{0,1,\infty\}=U_{0,3}(\mathbb C)$,
$G=\mathfrak S_3$ operating in the usual way, thus (with the \deleted{\unsure{notations}} \inserted{conventions} of the preceding \S),
\deleted{$Y=\unsure{\mathbb P^1(\mathbb C)}\,\mathbb P^1(\mathbb C)$}
\inserted{identifying} $\mathbb P^1(\mathbb C)/(G=\mathfrak S_3)$ with $\mathbb P^1_{\mathbb C}$, with
\[
(4)\qquad
\begin{array}{rcll}
(0,1,\infty)&\longmapsto&0&d_0=2\\
(2,-1,\tfrac12)&\longmapsto&1&d_1=2\\
(j,\overline j)&\longmapsto&\infty&d_\infty=3
\end{array}
\]
we find $Y=\mathbb P^1(\mathbb C)\setminus\{0\}$, and $\underline d=2\{1\}+3\{\infty\}$.

We thus find a canonical isomorphism
\[
(5)\qquad
\Pi^{1}((U_{0,3}(\mathbb C),\mathfrak S_3),x)
\simeq\Pi_1((\mathbb P^1(\mathbb C)\setminus\{0\},\underline d),\unsure{y})
\]

\scanpage{536}

\unsure{now|or} unfortunately we cannot take $x=P=j$, we have to take
$x\in\mathbb P^1(\mathbb C)\setminus\{0,1,\infty,2,-1,\tfrac12,j,\overline j\}$
--- the most convenient thing will be to take $x$ such that $y=f(x)\,\unsure{\in|\leq|\leftarrow}$ is the \inserted{typical} base-$\mathrm p^{\mathrm t}$ $Q=j$ for calculating the $\pi_1$ with ramification --- for such an $x$ we thus find a canonical isomorphism
\[
(5\,\mathrm{bis})\qquad
\Pi_1((U_{0,3}(\mathbb C),\mathfrak S_3),x)\simeq\Pi'_{0,3}
\]
and by choosing an \inserted{$\mathfrak S_3$-}class of paths from the \inserted{initial} base point $P=j$ on $\unsure{X_0|X}=U_{0,3}(\mathbb C)$ to $x$, whence a corresponding isomorphism of the first side of (5bis) \deleted{\illegible{1}} with $\Pi_{0,3}$, we find a canonical isomorphism
\[
(6)\qquad
\Pi_1((U_{0,3}(\mathbb C),\mathfrak S_3),P)\simeq\Pi'_{0,3}
\]
(which changes by an inner automorphism when one modifies the choice of an $\mathfrak S_3$-class of paths), whence the exact sequence
\[
(7)\qquad
1\longrightarrow\Pi_{0,3}\longrightarrow\Pi'_{0,3}\longrightarrow\mathfrak S_3\longrightarrow 1,
\]
which is none other than (9bis) of the preceding \S\, but interpreted here much more transparently as
\scanpage{537}
\mspage{274}

the homotopy exact sequence for $\mathfrak{S}_3$ operating
on $U_{0,3}(\mathbb{C})$. And we find likewise,
in the \unsure{arithmetico-geometric} setting,
working on the scheme $U_{0,3}=U_{0,3\mathbb{Q}}$ and its
quotient $U_{0,3}/\mathfrak{S}_3\simeq\mathbb{P}^{1}_{\mathbb{Q}}\setminus\{0\}$, with
the signature $2\{1\}+\unsure{3}\{\infty\}$ on it, the
diagram of exact sequences

\diagram{
\begin{tikzcd}[ampersand replacement=\&,column sep=large,row sep=large]
 \& 1 \arrow[d] \& 1 \arrow[d] \& \& \\
1 \arrow[r] \& \widehat{\Pi}_{0,3} \arrow[r] \arrow[d] \& \widehat{\Pi}'_{0,3} \arrow[r] \arrow[d] \& \mathfrak{S}_3 \arrow[r] \arrow[d,equal] \& 1 \\
1 \arrow[r] \& \mathcal{E}_{0,3} \arrow[r] \arrow[d] \& \mathcal{E}'_{0,3} \arrow[r] \arrow[d] \& \mathfrak{S}_3 \arrow[r] \& 1 \\
 \& \Gamma_{\mathbb{Q}} \arrow[r,equal] \arrow[d] \& \Gamma_{\mathbb{Q}} \arrow[d] \& \& \\
 \& 1 \& 1 \& \&
\end{tikzcd}
[(8) to the left of the second horizontal sequence.
To the left and a little below $\widehat{\Pi}_{0,3}$:
$\widehat{\mathcal{T}}^{+}_{0,4}$, followed by two short parallel oblique strokes
(an $=$ sign equating it with $\widehat{\Pi}_{0,3}$).
To the left and a little below $\mathcal{E}_{0,3}$:
$\unsure{N}^{!}_{0,4}$ (a short horizontal stroke against the ${}^{!}$), followed by two short
parallel oblique strokes (an $=$ sign toward $\mathcal{E}_{0,3}$).
To the left and a little below the first $\Gamma_{\mathbb{Q}}$:
$\Gamma_{0,4}$, followed by two short parallel oblique strokes (an $=$ sign toward $\Gamma_{\mathbb{Q}}$).]
}

We now want to relate
$U_{0,3}$, with the action of $\mathfrak{S}_3$
on it, to the modular multiplicity
$M_{11}$ for elliptic curves
(subscripts $\mathbb{Q}$ understood). It is
a pity, by the way, to work only over $\mathbb{Q}$ --- I am going to work
instead, over $\mathbb{Z}[1/2]$ --- i.e. in \unsure{residue|always} characteristic $\ne 2$.

\scanpage{538}

In the preceding \S{} I identified, somewhat
murkily, $U_{0,3}=M^{!}_{0,4}$\unsure{,} with $M_{1,1}[2]'$\unsure{,}
(\deleted{of which}\inserted{over which} $M_{1,1}[2]$
\deleted{\unsure{should be|was to be} \illegible{1}}\inserted{is a gerbe banded} by
$\mathbb{Z}/2$). The technical troubles (rather,
the conceptual complications) are due
to the fact that for elliptic curves
(more generally, for abelian varieties of
arbitrary dimension), if $n$ is an integer $\ge 3$,
the ``Jacobi rigidification of level $n$''
really is a rigidification (i.e.
every automorphism of a relative elliptic curve
$E$ which induces the identity on
${}_nE$ --- kernel subscheme of $n.\mathrm{id}_E$ --- is
the identity --- but it is no longer the
same for $n=2$, where $\unsure{u|n}$ is the \uline{identity}\inserted{$\mathrm{id}_E$} \unsure{up} above
an open and closed subset of $S$, and $-\mathrm{id}_E$
on the complement. This suggests, for
a better understanding, covering
$M_{1,1}[2]$ by an $M_{1,1}(n)$, where $n$ is a
multiple of $2$ --- then $M_{1,1}(n)$ becomes an
ordinary scheme of finite type over $\operatorname{Spec}\mathbb{Z}$
(its image in $\operatorname{Spec}\mathbb{Z}$ \unsure{will be|will be} \inserted{the open set} $\operatorname{Spec}\mathbb{Z}[1/n]$)
and over $\operatorname{Spec}[1/n]$, the \deleted{\unsure{scheme|scheme}}\inserted{modular topos}
\scanpage{539}
\mspage{275}

$M_{1,1}$ is recovered as
\[
(9)\qquad
M_{1,1}\unsure{\simeq|\cong}
\bigl(M_{1,1}[n],\overbrace{\mathrm{Gl}(2,\mathbb Z/n\mathbb Z)}^{\Gamma_n}\bigr)
\]
(over $\operatorname{Spec}\mathbb Z[1/n]$)

\margin{careful, it really is $\mathrm{Gl}(2,\mathbb Z/n\mathbb Z)$, and \uuline{not} $\mathrm{Sl}(2,\mathbb Z/n\mathbb Z)$ that one must take here}

Here $M_{1,1}[n]$ is the scheme representing the functor on $(\mathrm{Sch})_{\unsure{\mathbb Z|S}}$
\[
S\longmapsto
\begin{gathered}
\text{set of relative elliptic curves}\\
\text{over $S$, \uline{equipped} with an isomorphism}\\
(\mathbb Z/n\mathbb Z)^2_S\longrightarrow {}_nE
\end{gathered}
\]
\margin{this functor \unsure{is|$\to$} representable iff $n\geq 3$ (so it is not for $n=1,2$)}
on which the group $\Gamma_n=\mathrm{Gl}(2,\mathbb Z/n\mathbb Z)=\operatorname{Aut}((\mathbb Z/n\mathbb Z)^2)$ operates in the obvious way. \deleted{The variety} The \uline{modular scheme} (``coarse'', as opposed to the modular multiplicity or topos) is described over $\mathbb Z/n\mathbb Z$ by
\[
(10)\qquad \widetilde M_{1,1}\simeq M_{1,1}[n]/\Gamma_n
\]
(over $\mathbb Z[1/n]$)

\uline{N.B.} The scheme $\widetilde M_{1,1}$ \inserted{over $\operatorname{Spec}\mathbb Z$} can be described as ``the representable envelope'' of the nonrepresentable functor
\[
S\longmapsto
\begin{gathered}
\text{\unsure{isomorphism|Isomorphism} classes of relative}\\
\text{elliptic curves over $S$ \ldots}
\end{gathered}
\]
-- ditto over a base scheme (for example $\mathbb Q$, or $\operatorname{Spec}\mathbb Z[1/n]$, arbitrary \ldots

\scanpage{540}

One must be careful that as a scheme over $\mathbb Z[1/n]$ (or over $\mathbb Q$, passing to the \unsure{generic|generic} fiber) $M_{1,1}[n]$ \uline{is not} relatively connected (i.e. does not have geometrically connected fibers). \unsure{In} fact, an isomorphism
\[
(\mathbb Z/n\mathbb Z)^2_S\simeq {}_nE
\]
implies, by passing to the $2^{\unsure{\mathrm{ième}|\mathrm{ème}}}$ exterior power, an isomorphism
\[
(\mathbb Z/n\mathbb Z)_S\xrightarrow{\sim}
\Lambda^2_{\mathbb Z/n\mathbb Z}\,{}_nE
\simeq \mu_n^{\otimes-1}(S)
\]
whence an isomorphism
\[
\mathbb Z/n\mathbb Z\simeq\mu_n(S)
\]
which identifies (taking \deleted{\unsure{and}} the image of $1$ mod $n$) with a section of $\mu_n^*(S)$, where $\mu_n^*$ is the relative $(\mathbb Z/n\mathbb Z)^*$-torsor over $\operatorname{Spec}\mathbb Z[1/n]$ of ``primitive \unsure{$n$th|$n$-th} roots of $1$''. We thus have a canonical morphism
\[
(11)\qquad M_{1,1}[n]\longrightarrow\mu_n^*
\]
and it is this morphism that has geometrically connected fibers \inserted{\uuline{in characteristic 0}}. Passing to the limit over variable $n$, we find \uline{over $\mathbb Q$}
\[
(12)\qquad M_{1,1}[\infty]=\varprojlim_n M_{1,1}[n]
\longrightarrow\mu_\infty^*\overset{\inserted{\mathrm{iso.can}}}{\simeq}\operatorname{Spec}\Sigma
\]
